
\section{Parsing}\fbeg{}\begin{frame}{\ft{Parsing for an Embedded Compiler}}

\vspace{.5em}	

\begin{center}

%\MyDiamond{}
{\Large%\fontfamily{uhv}\selectfont
%\begin{center}
\hspace*{-2pt}\begin{minipage}{\textwidth}
\vspace{4pt}

%\definecolor{blback}{RGB}{0,100,100}	
%\definecolor{blfront}{RGB}{0,100,50}

%\fcolorbox{lqboutercolor}{lqbinnercolor}{\begin{minipage}{\textwidth}%
		


{\hspace{1.5em}\begin{minipage}[l]{.9\textwidth}\Large\centering\color{slidePartHeadColor} 	
{\color[rgb]{.3,.1,0}{\Huge\fontfamily{bch}\fontseries{eb}\selectfont Useful 
design patterns for parsers oriented to self-contained compiler 
stacks 

}}\\\vspace{1em}
{\LARGE \textbf{Avoid relying on external tools such as Lex, Yacc, Bison, LLVM, etc.}}
\vspace{.7em}	
\end{minipage}}
\vspace{.1em}

\begin{minipage}{.88\textwidth}
{\LARGE \setlength{\leftmargini}{30pt}\begin{enumerate}
\dmitem {\dksp} Grammars should be context-sensitive: they might 
have one or more notions of \q{context} to support language extensions and expressive syntax \vspace{10pt}

\dmitem {\dksp} Individual rules are regular expression objects coupled with callback 
handlers, and optionally contextual filters: we are not generating code 
\textsl{for} a parser, we are just implement rule-lists that are tried 
in sequence --- anchored 
at the current \q{cursor} position --- until one regex matches  \vspace{10pt}

\dmitem {\dksp} Rule callbacks can modify parser state/contexts, but in general 
would emit code for a \q{first-stage} Virtual Machine \vspace*{11pt}

\dmitem {\dksp} This preliminary VM then calls instructions for a code generator 
that produces \q{bytecode} for the main VM \vspace*{11pt}

\dmitem {\dksp} In both contexts \q{bytecode} is actually calls to class methods, 
and can be printed in human-readable fashion (although compressed bytecode can employ single bytes per instruction)


\end{enumerate}}
\end{minipage}

\end{minipage}

%}
%\end{minipage}
%\end{center}
}

\end{center}


\end{frame}

